%! AP Statistics — Unit 3, Lesson 3.7 — Homework KEY
\documentclass[10pt]{article}
\usepackage{apstats-article}
\usepackage{apstats-key}

\begin{document}

\pageheader{Unit 3, Lesson 3.7}{Homework — KEY}
\namedateperiod

\vspace{0.1in}

\begin{practicebox}
\small
\textbf{1.}

Rand. assignment: \ans{No} \quad Rand. selection: \ans{Yes}\\[3pt]
Causation: \ans{No} \quad Generalize to: \ans{All registered voters nationally}\\[3pt]
\ans{Because no random assignment was used, we cannot conclude that social media use causes more polarized views — confounding variables (e.g., age, news consumption) may explain the association. Because a random sample was used, the finding can be generalized to all national registered voters.}

\vspace{6pt}
\textbf{2.}

Rand. assignment: \ans{Yes} \quad Rand. selection: \ans{No (volunteers)}\\[3pt]
Causation: \ans{Yes} \quad Generalize to: \ans{Similar volunteers; not all students}\\[3pt]
\ans{Because random assignment was used, the active-learning classroom caused the significantly higher exam scores. Because the participants were volunteers (not randomly selected), we cannot generalize beyond students similar to these volunteers.}

\vspace{6pt}
\textbf{3.}

Rand. assignment: \ans{Yes} \quad Rand. selection: \ans{Yes}\\[3pt]
Causation: \ans{Yes} \quad Generalize to: \ans{Patients at these 10 hospitals (or similar hospital patients)}\\[3pt]
\ans{Because random assignment was used, the new drug caused the statistically significant improvement. Because patients were randomly selected from these hospitals, we can generalize the finding to the patient population at these hospitals.}

\vspace{6pt}
\textbf{4.} \ans{Statistical significance means the observed result is too large to be explained by chance alone. Practical importance (effect size) refers to whether the difference is large enough to matter in real life. Example: A study finds that a new drug reduces blood pressure by 1 mm Hg (statistically significant with n = 100,000). This is statistically significant but probably not clinically important — a 1 mm Hg difference has little real-world impact on health outcomes.}
\end{practicebox}

\end{document}
